\documentclass[11pt,a4paper]{article}

% ---------------------------------------------------------------------------
% Spatial Void Feature Extraction from X-ray Images
% Companion document to notebooks/jp_void_spatial_features.ipynb
% AI 221 Project -- Machine Learning-Based Classification of Thermal
% Resistance Degradation in Solder Joints Using Void Spatial Distribution
% Features from X-ray Images
% ---------------------------------------------------------------------------

\usepackage[utf8]{inputenc}
\usepackage[T1]{fontenc}
\usepackage{lmodern}
\usepackage[margin=2.3cm]{geometry}
\usepackage{amsmath,amssymb}
\usepackage{booktabs}
\usepackage{longtable}
\usepackage{array}
\usepackage{tabularx}
\usepackage{ragged2e}
\usepackage{enumitem}
\usepackage{graphicx}
\usepackage{xcolor}
\usepackage{listings}
\usepackage{url}
\usepackage[hidelinks]{hyperref}
\usepackage{caption}
\usepackage{float}

\captionsetup{font=small,labelfont=bf}
\setlength{\parskip}{0.4em}
\setlength{\parindent}{0pt}

% column types for the long tables
\newcolumntype{L}[1]{>{\RaggedRight\arraybackslash}p{#1}}

% inline code style
\definecolor{codegray}{gray}{0.95}
\lstset{
  basicstyle=\ttfamily\small,
  backgroundcolor=\color{codegray},
  breaklines=true,
  frame=single,
  framerule=0pt,
  xleftmargin=0.5em,
  language=Python,
  showstringspaces=false,
  keywordstyle=\bfseries,
  commentstyle=\itshape\color{gray}
}
% code tokens may break after an underscore so long names fit in narrow table columns
\newcommand{\code}[1]{{\ttfamily\def\_{\textunderscore\allowbreak}#1}}

\title{Spatial Void Feature Extraction from X-ray Images\texorpdfstring{\\
\large Definitions, Computation and Relevance to Thermal Resistance Degradation}{: Definitions, Computation and Relevance to Thermal Resistance Degradation}}
\author{AI 221 Project \\ Jocelyn de Castro, Joanna Marie Pagtalunan}
\date{November 2026}

\begin{document}
\maketitle

\begin{abstract}
This document explains the spatial void feature extraction pipeline implemented in the
project notebook \code{jp\_void\_spatial\_features.ipynb}. It covers how the X-ray images of the
\emph{hellastudy} LED dataset are registered, how a void-free reference image is built for each
LED type, how solder pads are located, how voids are segmented, and how the final per-LED
feature vector is assembled. For every extracted feature we give a definition, the formula or
procedure used to compute it, its range, and the reason it is expected to carry information
about thermal resistance degradation, drawing on the project proposal and the related
literature in the \code{docs/} folder.
\end{abstract}

\tableofcontents

% ===========================================================================
\section{Motivation}
% ===========================================================================

The project proposal \cite{proposal} observes that industry acceptance criteria for solder
voids (IPC-A-610, IPC-7095, IPC-7093, MIL-STD-883) are based almost entirely on the
\emph{total} void area, typically a limit around 25\,\%. None of them relates the
\emph{position} of a void to thermal performance. The proposal argues that a void of a given
size located directly beneath the die on the thermal pad obstructs the primary heat path and
should degrade thermal resistance more than an equally large void at the periphery of an
electrical pad, yet both pass or fail the same way under current rules.

The hellastudy dataset \cite{part1,part2} provides what is needed to test this: an X-ray image of
every one of the 1\,800 LED solder joints at zero thermal shock cycles (TSC), plus transient
thermal analysis (TTA) measurements of thermal resistance $R_{th}$ throughout 1\,500 cycles of
thermal shock. Earlier work on this dataset predicted thermal properties from scanning acoustic
microscopy (SAM) images using convolutional networks \cite{zippelius2022}, and later showed that
SAM images add spatial information that TTA alone lacks and improve long-range reliability
prediction \cite{zippelius2026}. Those studies let the network learn its own image features. Our
approach is complementary: we extract a small set of \emph{interpretable} location-aware void
descriptors from the X-ray images so that the contribution of void position can be examined
directly and fed to classical learners (Random Forest, XGBoost, SVM) as well as an ANN.

The notebook therefore has two goals:
\begin{enumerate}[nosep]
  \item Reproduce the void segmentation of Part~I of the hellastudy \cite{part1} closely enough
        that the total void ratio matches the authors' own published per-LED values.
  \item Go beyond the void ratio and compute where the voids sit on each pad.
\end{enumerate}

% ===========================================================================
\section{Data}
% ===========================================================================

The X-ray images live in one folder per LED type under \path{data/raw/XRay/XRay/}, with file
names of the form
\begin{center}
\path{XRay_<LED_TYPE>_<PASTE>_Panel<N>_LED<NN>_<TSC>TSC.jpg}
\end{center}
The regular expression in \code{parse\_filename} splits this into \code{led\_type}, \code{paste}, \code{panel},
\code{led\_num} and \code{tsc}, and builds the key
\code{unit\_id = <LED\_TYPE>\_<PASTE>\_P<panel>\_L<led\_num>} that is later used to join the
features with the thermal-resistance target.

\begin{itemize}[nosep]
  \item 9 LED types $\times$ 5 solder pastes $\times$ 40 LEDs $=$ 1\,800 images, all at 0~TSC.
  \item Image sizes range from $200\times200$ to $300\times300$ pixels and are constant
        within an LED type.
  \item Six files are placeholder pictures that only contain the word ``Missing''. They are
        detected by \code{is\_placeholder} (more than 60\,\% of pixels at grey level $\geq 250$)
        and excluded from the reference image; their feature rows are set to \code{NaN}.
  \item Images are loaded as 8-bit greyscale. In X-ray, metal attenuates radiation, so pads
        and solder appear \emph{dark} while voids appear \emph{bright} (Part~I,
        \S III.A: ``voids appear brighter in the X-ray image since more radiation can pass
        through the specimen due to missing metallic material'' \cite{part1}).
\end{itemize}

Pad counts per LED type are taken from Part~I \S II.A and Part~II Table~I \cite{part1,part2}:
FC-SP1, FC-SP2 and FC-GB3 are symmetrical two-pad designs with no dedicated thermal pad; all
other types have three pads, of which the largest is the thermal pad (Part~II \S II.A: ``for
three-pad designs, the thermal pad is typically bigger compared to the electrical pads'').

% ===========================================================================
\section{Pipeline overview}
% ===========================================================================

Part~I of the hellastudy describes its void algorithm as follows (\S IV.A, Fig.~5): the outline
of all LEDs of one type is detected, an average image is generated by overlapping all of them to
obtain a ``void-less'' X-ray image, this average is subtracted from each device under test
(DUT), and the result is thresholded \cite{part1}. The notebook follows the same idea with a
more robust implementation, in six stages:

\begin{enumerate}
  \item \textbf{Grouping.} Images are processed one LED type at a time, because pad layout,
        magnification and image size differ between types.
  \item \textbf{Registration and reference image} (\S\ref{sec:registration}). All images of a
        type are aligned to a common frame and their per-pixel \emph{median} is taken as the
        void-free reference.
  \item \textbf{Pad detection} (\S\ref{sec:pads}). The pads are found automatically as the
        darkest connected components of the reference and assigned a role
        (thermal or electrical).
  \item \textbf{Void mask} (\S\ref{sec:voidmask}). The reference is fitted locally to each DUT
        and subtracted. The residual is thresholded with hysteresis to produce a binary void
        mask.
  \item \textbf{Validation} (\S\ref{sec:validation}). The resulting void ratio of every LED is
        compared with the authors' published values in
        \code{data/raw/CrackVoid Ratios/Xray Void Ratio.csv}.
  \item \textbf{Feature extraction} (\S\ref{sec:features}). Spatial features are computed per pad
        in pad-normalised coordinates and aggregated per LED by pad role.
\end{enumerate}

The output is \code{data/processed/spatial\_void\_features.csv}, a matrix of 1\,799 rows (one
LED is absent from the FC-GB2 folder) and 32 columns.

% ===========================================================================
\section{Registration and reference image}
\label{sec:registration}
% ===========================================================================

A plain average of unaligned images blurs pad edges and internal structures, and those blurred
edges then survive subtraction and get counted as voids. The notebook therefore registers each
image before building the reference, in \code{register\_image} and \code{register\_led\_type}.

\paragraph{Per image.}
\begin{enumerate}[nosep]
  \item \emph{Coarse translation search.} The central 76\,\% crop of the template is searched
        over the DUT (padded by 60~px) with normalised cross-correlation
        (\code{cv2.matchTemplate}, \code{TM\_CCOEFF\_NORMED}). This handles LEDs that sit tens of
        pixels off-centre. The correlation score is kept as a quality indicator.
  \item \emph{ECC refinement, Euclidean.} \code{cv2.findTransformECC} with
        \code{MOTION\_EUCLIDEAN} refines the shift and adds rotation.
  \item \emph{ECC refinement, affine (last pass only).} A second ECC run with
        \code{MOTION\_AFFINE} absorbs the $\pm1.5\,\%$ magnification differences between shots.
        The affine result is accepted only if both scale factors lie in $(0.94, 1.06)$.
  \item The DUT is warped into the template frame. A \emph{validity mask} is produced by
        warping an all-ones image the same way; pixels that fall outside the original frame are
        marked invalid and excluded from all later statistics.
\end{enumerate}

\paragraph{Per LED type.} Three passes are run. Pass~1 aligns every image to the first real
image of the type. Passes~2 and~3 align to the per-pixel \emph{median} of the previous pass
(restricted to non-placeholder images with correlation score $> 0.5$). The median is used
instead of the mean because voids are at random positions in different LEDs and the median
suppresses them, whereas the mean leaves faint ghosts. The reference image is the median of the
final pass. Registration results are cached in \code{\_REG\_CACHE}.

% ===========================================================================
\section{Pad detection}
\label{sec:pads}
% ===========================================================================

Pads are located once per LED type on the reference image by \code{detect\_pads}. The physical
basis is again X-ray attenuation: pad plus solder is the thickest metal stack in the image and
therefore its darkest large region.

\begin{enumerate}[nosep]
  \item \emph{Dark-pixel mask.} By default Otsu's threshold (inverted). Several LED types use
        a percentile override from \code{PAD\_THRESH\_OVERRIDE}, e.g.\ the darkest 40\,\% of
        pixels for FC-SP4, FC-GB1 and VTF2, because the die or a bump grid competes with the pads
        in darkness.
  \item \emph{Iterative erosion.} The mask is eroded with a $3\times3$ kernel until at least
        the expected number of components (2 or 3 from \code{PAD\_COUNT}) with area
        $\geq 2\,\%$ of the image exist, up to \code{MAX\_EROSION\_OVERRIDE} passes (8, or 20
        for VTF2). This separates pads that touch after thresholding without needing a visible
        gap.
  \item \emph{Keep the $N$ largest components.}
  \item \emph{Dilate back and clip.} Each component is dilated by the same number of erosions
        and intersected with the original dark-pixel mask, which recovers the full pad extent
        without bleeding into a neighbour.
  \item \emph{Role assignment.} For three-pad types the largest component is the
        \emph{thermal} pad and the others are \emph{electrical}; for two-pad types both are
        electrical.
\end{enumerate}

Each pad is stored as a bounding box $(x_0, y_0, x_1, y_1)$, a role and a pixel area, together
with its exact pixel mask. For the FC-SP3 example in the notebook the thermal pad spans
$156\times88$ pixels (13\,500 pad pixels) and the two electrical pads about $69\times88$
pixels each.

% ===========================================================================
\section{Void mask extraction}
\label{sec:voidmask}
% ===========================================================================

A plain difference $\mathrm{DUT} - \mathrm{reference}$ is not sufficient on these images: exposure
and blur differ between shots, varying solder thickness adds slow brightness gradients, and the
die, bumps and wire bonds sit in slightly different places relative to the pads. In the
notebook's diagnostic cell, the plain difference inside pads has mean 1.6 and 99th percentile
15 grey levels, which is of the same order as the void signal. The notebook therefore fits the
reference to each DUT \emph{locally} before subtracting (\code{\_local\_residual}).

\subsection{Local linear model}

Within a Gaussian window of $\sigma_w = 8$~px around every pixel, the DUT is modelled as
\begin{equation}
  \mathrm{DUT}(x,y) \approx b + a\,\mathrm{ref} + c\,\frac{\partial\,\mathrm{ref}}{\partial x}
  + d\,\frac{\partial\,\mathrm{ref}}{\partial y} + e\,\nabla^2 \mathrm{ref},
  \label{eq:localfit}
\end{equation}
where the five regressors are built by \code{\_ref\_features} (the reference, its Sobel
derivatives and its Laplacian, after a Gaussian blur of $\sigma = 1$). The coefficients have a
physical reading:
\begin{itemize}[nosep]
  \item $a, b$: local contrast and brightness (exposure differences, solder-thickness gradients);
  \item $c, d$: a local sub-pixel shift (first-order Taylor expansion of a translation);
  \item $e$: a blur difference (first-order expansion of a Gaussian blur).
\end{itemize}
The normal equations are solved per pixel using Gaussian-weighted moments, with a small ridge
penalty ($\lambda = 4$) that shrinks the contrast coefficient towards the \emph{global} gain
$a_0$ estimated robustly by \code{\_global\_gain} (iteratively reweighted linear fit with a
$2.5\sigma$ MAD cut). The fit is iterated three times; after each iteration, pixels whose
residual exceeds $2\sigma$ (bright) or $-4\sigma$ (dark), with $\sigma$ the MAD-based robust
scale, are removed from the fit weights. Because voids are \emph{not} present in the reference
they remain in the residual, and because they are removed from the fit they do not pull the
local brightness estimate upward.

The residual is divided by $\max(a_0, 0.05)$ so that it is expressed in grey levels of the
\emph{reference}, making the detection threshold comparable across shots with different
exposure.

\subsection{Regions}

\code{build\_void\_context} prepares, once per LED type:
\begin{itemize}[nosep]
  \item \code{pad}: union of all pad masks;
  \item \code{fit\_region}: pad dilated by 4~px (area over which the local fit is computed);
  \item \code{detect\_region}: pad eroded by 2~px (voids are only accepted here, to stay away
        from the pad outline where registration offsets concentrate);
  \item \code{dark}: small very dark package features (bump grid, vias, wire-bond balls) found
        by a black-hat transform with an $11\times11$ ellipse and dilated by 2~px. These are
        never solder and are excluded from both fit and detection;
  \item \code{edge\_dist}: Euclidean distance transform of the detect region, used to recognise
        blobs hugging the pad outline.
\end{itemize}

\subsection{Thresholding and blob filtering}

\code{extract\_void\_mask} applies a hysteresis threshold to the residual: a blob must exceed
$T$ somewhere and is then grown down to $0.6\,T$. $T$ is set per LED type in \code{VOID\_T}
(10 to 16 grey levels). Blobs are then filtered:
\begin{itemize}[nosep]
  \item area $< 8$~px are dropped (the reference CSV contains no void smaller than 8~px);
  \item elongation (major/minor axis) $\geq 3$ with at least 60\,\% of pixels within 3~px of the
        pad edge are dropped (pad or solder outline offsets);
  \item elongation $\geq 4$ with major axis $\geq 30$~px are dropped (edges of shifted internal
        structures).
\end{itemize}
What survives is the binary void mask for that DUT.

\subsection{Calibration of \texorpdfstring{$T$}{T}}

\code{calibrate\_threshold} sweeps $T$ from 6 to 20 for one LED type and compares the resulting
per-LED void ratio with the CSV. For FC-SP3 the mean absolute error is minimised at $T = 12$
(0.38 percentage points, Pearson $r = 0.92$), which is the value stored in \code{VOID\_T}. The
same procedure was used for every type.

% ===========================================================================
\section{Validation against the published void ratios}
\label{sec:validation}
% ===========================================================================

\code{Xray Void Ratio.csv} holds the hellastudy authors' own result for every LED (pad pixels,
void count, void pixels, void ratio); its per-type means reproduce Table~III of Part~I
\cite{part1}. One duplicated key (FC-SP4 / SAC105 / panel~2 / LED~7) is dropped as ambiguous.
Table~\ref{tab:validation} compares the notebook's void ratio with the CSV for every LED type.

\begin{table}[H]
\centering
\caption{Void ratio from this notebook versus the authors' CSV, per LED type.
Errors are in percentage points of void ratio.}
\label{tab:validation}
\small
\begin{tabular}{lrrrrrr}
\toprule
LED type & $n$ & Ours mean [\%] & CSV mean [\%] & MAE [pt] & Pearson $r$ & Spearman $\rho$ \\
\midrule
FC-GB1 & 200 & 2.37 & 2.35 & 0.48 & 0.90 & 0.92 \\
FC-GB2 & 199 & 1.31 & 1.34 & 0.24 & 0.95 & 0.92 \\
FC-GB3 & 198 & 4.96 & 5.33 & 1.08 & 0.92 & 0.93 \\
FC-SP1 & 200 & 0.45 & 0.51 & 0.25 & 0.49 & 0.71 \\
FC-SP2 & 198 & 2.25 & 2.45 & 0.57 & 0.84 & 0.92 \\
FC-SP3 & 200 & 2.33 & 2.37 & 0.38 & 0.92 & 0.94 \\
FC-SP4 & 198 & 3.18 & 3.17 & 0.48 & 0.89 & 0.95 \\
VTF1   & 198 & 2.58 & 2.86 & 0.75 & 0.79 & 0.82 \\
VTF2   & 200 & 3.39 & 3.47 & 0.81 & 0.83 & 0.81 \\
\midrule
All    & 1\,791 & \multicolumn{2}{c}{} & 0.56 & 0.91 & \\
\bottomrule
\end{tabular}
\end{table}

The per-type means agree with Part~I Table~III to within a few tenths of a percentage point.
FC-SP1 has the lowest correlation because its voiding is so small (mean 0.5\,\%) that a few
pixels dominate the ratio. This validation matters because it shows the void masks on which all
spatial features are built are consistent with the segmentation used by the original authors,
so any predictive power we find for the spatial features is not an artefact of a different
segmentation.

% ===========================================================================
\section{Spatial feature extraction}
\label{sec:features}
% ===========================================================================

\subsection{Coordinate system}

All features are computed by \code{pad\_features} on the void mask cropped to the pad bounding
box, in \emph{pad-normalised coordinates}: the pad centre is $(0,0)$ and the corners are
$(\pm1,\pm1)$. For a pad of width $w$ and height $h$ pixels and a void pixel at image column $x$
and row $y$ within the crop,
\begin{equation}
  \tilde{x} = \frac{x - w/2}{w/2}, \qquad \tilde{y} = \frac{y - h/2}{h/2}.
\end{equation}
This normalisation makes features comparable across LED types whose pads range from
$0.3$~mm$^2$ to $2.7$~mm$^2$ \cite{part2}. Note that the image $y$ axis points \emph{down}, so
$\tilde{y} < 0$ means the upper half of the pad as seen in the image, and the quadrant labels
NW, NE, SW, SE refer to the image frame (N = top of the image). For a three-pad LED imaged as
in the notebook figures, the thermal pad is above the two electrical pads, so ``S'' on the
thermal pad is the side facing the electrical pads and the inter-pad gap.

\subsection{Per-pad features}

Table~\ref{tab:padfeatures} defines the eleven features computed for each pad. In the formulas,
$V$ is the set of void pixels inside the pad crop, $|V|$ its cardinality, $A_{\mathrm{pad}} = w\,h$
the bounding-box area, and $\{R_k\}$ the connected components (8-connectivity, via
\code{skimage.measure.label}) of the void mask inside the crop.

\begin{longtable}{L{2.8cm} L{4.4cm} L{5.2cm} L{2.3cm}}
\caption{Per-pad spatial void features computed by \code{pad\_features}.}
\label{tab:padfeatures}\\
\toprule
\textbf{Feature} & \textbf{Definition} & \textbf{How the value is obtained} & \textbf{Range} \\
\midrule
\endfirsthead
\toprule
\textbf{Feature} & \textbf{Definition} & \textbf{How the value is obtained} & \textbf{Range} \\
\midrule
\endhead
\bottomrule
\endfoot

\code{void\_area\_frac} &
Fraction of the pad area occupied by voids. &
$|V| / A_{\mathrm{pad}}$, i.e.\ \code{pad\_crop.sum() / (pad\_h * pad\_w)}. Uses the
bounding-box area as denominator; for the near-rectangular pads in this dataset this is close
to the exact pad-mask area used for \code{void\_ratio}. &
$[0, 1]$ \\

\code{centroid\_x} &
Horizontal position of the void centre of mass, relative to the pad centre. &
Mean column of all void pixels, normalised: $(\bar{x} - w/2)/(w/2)$. $-1$ = left edge, $0$ =
centre, $+1$ = right edge. Set to 0 when the pad has no voids. &
$[-1, 1]$ \\

\code{centroid\_y} &
Vertical position of the void centre of mass, relative to the pad centre. &
Mean row of all void pixels, normalised: $(\bar{y} - h/2)/(h/2)$. $-1$ = top edge (image
frame), $+1$ = bottom edge. Set to 0 when there are no voids. &
$[-1, 1]$ \\

\code{q\_NW}, \code{q\_NE}, \code{q\_SW}, \code{q\_SE} &
Share of the void pixels that lie in each quadrant of the pad. &
Pad is split at $x = w/2$ and $y = h/2$. Each value is (void pixels in that quadrant)
$/\,|V|$; the four values sum to 1. All 0 when there are no voids. &
$[0, 1]$ each \\

\code{cluster\_count} &
Number of separate voids on the pad. &
Number of connected components $\{R_k\}$ of the void mask inside the crop. &
$\{0, 1, 2, \dots\}$ \\

\code{mean\_circularity} &
Average roundness of the voids. &
For each component, circularity $= 4\pi A_k / P_k^2$ with $A_k$ the area and $P_k$ the
perimeter from \code{regionprops}; then the mean over components. A perfect disc gives 1,
elongated or ragged shapes give smaller values (pixelisation can push it slightly above 1 for
tiny blobs). \code{NaN} when there are no voids. &
$(0, \approx 1]$ or \code{NaN} \\

\code{largest\_void\_frac} &
Area of the single biggest void as a fraction of the pad. &
$\max_k A_k / A_{\mathrm{pad}}$. Equals \code{void\_area\_frac} when there is only one void. &
$[0, 1]$ \\

\code{edge\_proximity\_min} &
How close the nearest void comes to the pad edge, normalised so that 1 means the pad centre. &
A Euclidean distance transform is computed from the crop border inward
(\code{ndi.distance\_transform\_edt}); the minimum of that distance over all void pixels is
divided by $\min(w, h)/2$. 0 = a void touches the border, 1 = the closest void pixel is at the
pad centre. \code{NaN} when there are no voids. &
$[0, 1]$ or \code{NaN} \\
\end{longtable}

\subsection{Per-LED aggregation}

\code{led\_feature\_vector} turns the per-pad dictionaries into one row per LED. The thermal pad
is kept separately because it is the primary heat path; the electrical pads (one or two) are
averaged because the dataset does not identify anode and cathode. Table~\ref{tab:ledcolumns}
lists every column of the final matrix.

\begin{longtable}{L{4.6cm} L{10.6cm}}
\caption{Columns of \code{spatial\_void\_features.csv} (one row per LED, keyed by \code{unit\_id}).}
\label{tab:ledcolumns}\\
\toprule
\textbf{Column} & \textbf{Meaning} \\
\midrule
\endfirsthead
\toprule
\textbf{Column} & \textbf{Meaning} \\
\midrule
\endhead
\bottomrule
\endfoot
\code{led\_type}, \code{paste}, \code{panel}, \code{led\_num} &
Metadata parsed from the file name. \code{led\_type} and \code{paste} are the categorical inputs
that the related work found important \cite{zippelius2026}. \\
\code{has\_thermal\_pad} &
1 for three-pad designs (FC-SP3, FC-SP4, FC-GB1, FC-GB2, VTF1, VTF2), 0 for two-pad designs
(FC-SP1, FC-SP2, FC-GB3). \\
\code{thermal\_<feature>} (11 columns) &
The eleven per-pad features of Table~\ref{tab:padfeatures} evaluated on the thermal pad.
\code{NaN} for two-pad LEDs. \\
\code{electrical\_mean\_<feature>} (11 columns) &
The eleven per-pad features averaged over all electrical pads of the LED, ignoring \code{NaN}
entries (e.g.\ circularity of a void-free pad). \\
\code{pad\_px} &
Number of valid pad pixels (union of all pad masks, intersected with the registration
validity mask). Same definition as the CSV column of the original authors. \\
\code{void\_count} &
Total number of connected void components on the whole LED. \\
\code{void\_px} &
Total number of void pixels on the whole LED. \\
\code{void\_ratio} &
\code{void\_px / pad\_px}. This is the conventional total void ratio and the quantity validated
in Table~\ref{tab:validation}. It plays the role of the \emph{baseline} feature in the project. \\
\code{image\_missing} &
1 if the X-ray file was a ``Missing'' placeholder; all feature columns are then \code{NaN}. \\
\end{longtable}

Column count: 4 metadata $+$ 1 \code{has\_thermal\_pad} $+$ 11 thermal $+$ 11 electrical $+$
4 whole-LED statistics $+$ 1 \code{image\_missing} $=$ 32.

% ===========================================================================
\section{Why these features are relevant to thermal degradation}
\label{sec:relevance}
% ===========================================================================

The physical argument throughout the literature is that heat generated at the LED junction
flows through the die, the submount, the solder joint and the substrate to the heat sink. Any
loss of metallic cross-section in the solder reduces the area available for conduction and so
raises the thermal resistance $R_{th}$ \cite{part1,zippelius2022}. Voids remove cross-section
at time zero; cracks remove it progressively during thermal cycling. The features in
Table~\ref{tab:padfeatures} describe \emph{how much} cross-section is lost and \emph{where}.
Table~\ref{tab:relevance} connects each feature to a specific finding in the proposal or the
reference papers.

\begin{longtable}{L{3.4cm} L{11.8cm}}
\caption{Relevance of each feature to thermal resistance degradation, with supporting evidence.}
\label{tab:relevance}\\
\toprule
\textbf{Feature(s)} & \textbf{Why it should matter} \\
\midrule
\endfirsthead
\toprule
\textbf{Feature(s)} & \textbf{Why it should matter} \\
\midrule
\endhead
\bottomrule
\endfoot

\code{void\_ratio}, \code{void\_area\_frac} &
The conventional metric and the only one used by IPC/MIL acceptance criteria \cite{proposal}.
Part~I found a measurable linear increase of the solder-and-substrate resistance $R_{th\text{-}S\&S}$
with void ratio for the two LED types with the smallest pad areas, FC-SP1 and FC-GB3
(Table~VI of \cite{part1}: $+0.030$ and $+0.016$~K/W per percent of voiding), and the authors
explain it by the reduction of conducting cross-section. For the other types the effect was
within noise at the low (under 5\,\%) voiding achieved by vacuum reflow. This is exactly why the
proposal argues that \emph{total} area is an insufficient descriptor and motivates the baseline
versus spatial comparison in the project plan: the spatial model must beat a model that already
knows the void ratio, LED type and paste. \\

\code{thermal\_*} vs.\ \code{electrical\_mean\_*} split, \code{has\_thermal\_pad} &
The thermal pad sits directly below the die and is the primary heat path for three-pad designs
(Part~I \S II.A: an ``electrical isolated pad below the die for improved heat management''
\cite{part1}). Part~II reports from simulation that ``cracks in the thermal pad have higher
thermal impact than cracks in the electrical pads'' \cite{part2}; the same geometric argument
applies to voids. Keeping thermal-pad features separate lets a model weight them differently
from electrical-pad features. Two-pad LEDs route all heat through their electrical pads, so
\code{has\_thermal\_pad} tells the model which regime applies. Part~II also found that pad
design (pad area, inter-pad gap, thermal-to-electrical pad ratio) drives lifetime, so pad role
is a meaningful axis. \\

\code{centroid\_x}, \code{centroid\_y} &
First-order summary of \emph{where} the void mass is. The proposal's central hypothesis is that
a void under the heat source is worse than the same void at the periphery \cite{proposal}.
TTA theory in \cite{zippelius2022,zippelius2026} notes that the position of a defect along the
thermal path changes when it appears in the $Z_{th}(t)$ curve; laterally, a void centred under
the die intercepts the densest heat flux. A centroid near $(0,0)$ on the thermal pad therefore
flags the geometry the proposal is most concerned about. \\

\code{q\_NW}, \code{q\_NE}, \code{q\_SW}, \code{q\_SE} &
A coarse spatial histogram that captures asymmetry the centroid alone cannot (two voids in
opposite corners have a central centroid but zero central mass). SAM images of the aged joints
show that cracking under thermal shock ``proceeds from the corners inwards''
\cite{zippelius2026}, and Part~II reports that ``intense cracking in a single solder pad has
higher thermal impact than a homogeneous crack distribution'' \cite{part2}. Quadrant shares show whether voids are concentrated in one
corner, where they would merge early with an advancing crack front, or spread evenly. For the
thermal pad the southern quadrants face the inter-pad gap, so the split also encodes
orientation relative to the other pads. \\

\code{largest\_void\_frac} &
Part~I states that ``too high void ratios in one interconnect, especially single big voids, lead
to an earlier separation because the heat flow inside the LED package is changed'' \cite{part1}.
IPC-7095 and IPC-7093 likewise address the size of an individual void in addition to the total
\cite{proposal}. A single large void blocks a contiguous area of the heat path, while the same
area split into many small voids leaves conduction paths between them. \\

\code{cluster\_count} &
Complements \code{largest\_void\_frac}: for the same total area, many small voids versus one
large void correspond to the homogeneous versus concentrated distributions that Part~II found to
have different thermal impact \cite{part2}. Void count is also one of the columns the hellastudy
authors recorded in their own CSV, so it is part of the established description of the joint. \\

\code{mean\_circularity} &
Gas voids formed during reflow are approximately circular. Irregular or elongated bright
regions in X-ray are more likely non-wetted areas or contamination, which Part~I notes are also
detected by X-ray and were suspected for FC-GB3's high voiding (surface-finish contamination)
\cite{part1}. Non-wetting is an interface defect that both reduces heat transfer and weakens
the mechanical bond, so shape carries information about the defect mechanism, not just its
area. \\

\code{edge\_proximity\_min} &
Thermo-mechanical cracks initiate at pad corners and edges where the CTE-mismatch stress is
highest and grow inward \cite{part2,zippelius2026}. A void that already touches the pad edge
shortens the remaining ligament a crack must traverse and removes anchoring at the most stressed
location, so it may accelerate the $R_{th}$ rise during cycling. Conversely, the proposal's
concern is a void in the centre; a high \code{edge\_proximity\_min} on the thermal pad indicates
that scenario. The feature therefore spans both failure narratives. \\
\end{longtable}

\paragraph{Connection to the related machine-learning work.}
Zippelius et al.\ \cite{zippelius2022} showed that a CNN can translate SAM images of these same
LEDs into TTA quantities ($R_{th}$, $B_{max}$, $z_{max}$), establishing that inspection images
contain information about thermal state. Their 2026 follow-up \cite{zippelius2026} found that
adding SAM images to TTA features gives a clear accuracy gain for pass/fail prediction of
$R_{th}$ degradation, and attributes part of this to the image encoding pad shape and location
information that TTA cannot provide. Both studies learn image features implicitly. Our explicit
features make the same spatial information available to models that cannot consume images
(Random Forest, XGBoost, SVM) and allow feature-importance analysis to answer the proposal's
question directly: does void \emph{location} add predictive value beyond void \emph{amount}?

% ===========================================================================
\section{Limitations and practical notes}
% ===========================================================================

\begin{itemize}
  \item \textbf{Low voiding.} Because the study used vacuum reflow, the average void ratio is
        below 4\,\% for all types except FC-GB3 \cite{part1}, and Part~II found no
        correlation between void ratio and lifetime at these levels \cite{part2}. Any signal
        from spatial features will be subtle; the thermal-pad features of high-voiding LEDs are
        where it is most likely to appear.
  \item \textbf{Missing values.} \code{mean\_circularity} and \code{edge\_proximity\_min} are
        \code{NaN} for void-free pads, and all \code{thermal\_*} columns are \code{NaN} for
        two-pad LEDs. These must be imputed or handled explicitly before PCA, UMAP, K-Means or
        SVM, which do not accept \code{NaN}. Tree models can use an indicator plus a sentinel.
  \item \textbf{Denominator.} \code{void\_area\_frac} divides by the bounding-box area while
        \code{void\_ratio} divides by the exact pad-mask pixel count; they differ slightly for
        pads whose mask is not a full rectangle.
  \item \textbf{Orientation.} Quadrant and centroid signs are in the image frame. They are
        consistent within an LED type (all images are registered to one frame) but the physical
        meaning of ``north'' may differ between types depending on how the panel was imaged.
  \item \textbf{Electrical pads are averaged.} If anode and cathode behave differently, the
        averaging hides it. The notebook's next-steps section suggests replacing
        \code{electrical\_mean\_*} with position-indexed per-pad columns if that matters.
  \item \textbf{Threshold sensitivity.} \code{VOID\_T} is calibrated against the authors' CSV.
        If registration or the local fit is changed, \code{calibrate\_threshold} must be re-run.
  \item \textbf{X-ray only.} SAM is deliberately not used for initial void detection because
        Part~I found it less sensitive than X-ray for small voids \cite{part1}; it becomes
        relevant for crack tracking during ageing, which is outside this notebook.
\end{itemize}

% ===========================================================================
\section{Next steps in the project}
% ===========================================================================

\begin{enumerate}[nosep]
  \item Join the feature matrix with \code{rth\_start\_end.csv} on \code{unit\_id}, derive the
        cycles-to-failure using Part~II's criterion of a 20\,\% relative increase in $R_{th}$
        \cite{part2}, and discretise into the degradation classes.
  \item Fit a baseline model on \code{led\_type}, \code{paste} and \code{void\_ratio}, then a
        spatial model that adds all \code{thermal\_*} and \code{electrical\_mean\_*} columns,
        and compare cross-validated per-class F1.
  \item Use PCA/UMAP and feature importances to identify which spatial descriptors drive any
        improvement, which is the direct test of the hypothesis in the proposal.
\end{enumerate}

% ===========================================================================
\begin{thebibliography}{9}

\bibitem{proposal}
J.~de Castro and J.~M.~Pagtalunan, ``AI221 Project Proposal: Machine Learning-Based
Classification of Thermal Resistance Degradation in Solder Joints Using Void Spatial
Distribution Features from X-ray Images,'' November 2026.
(\code{docs/AI\_221\_Project\_Proposal.pdf})

\bibitem{part1}
M.~Schmid, A.~Zippelius, A.~Han\ss{}, S.~B\"ockhorst, and G.~Elger, ``Investigations on
High-Power LEDs and Solder Interconnects in Automotive Application: Part~I---Initial
Characterization,'' \emph{IEEE Trans.\ Device Mater.\ Rel.}, vol.~22, no.~2, pp.~175--186,
June 2022.

\bibitem{part2}
M.~Schmid, A.~Zippelius, A.~Han\ss{}, S.~B\"ockhorst, and G.~Elger, ``Investigations on
High-Power LEDs and Solder Interconnects in Automotive Application: Part~II---Reliability,''
\emph{IEEE Trans.\ Device Mater.\ Rel.}, vol.~23, no.~3, pp.~419--429, Sept.\ 2023.

\bibitem{zippelius2022}
A.~Zippelius, J.~Hermann, T.~Strobl, A.~Hoffmann, M.~Schmid, and G.~Elger, ``Predicting thermal
resistance of solder joints based on Scanning Acoustic Microscopy using Artificial Neural
Networks,'' in \emph{Proc.\ IEEE 9th Electronics System-Integration Technology Conference
(ESTC)}, 2022, pp.~566--575.

\bibitem{zippelius2026}
A.~Zippelius, Z.~Mohd, M.~Schmid, and G.~Elger, ``Comparison of Different Input Data for the
Prediction of LED Solder Joints Using Artificial Neural Networks,'' \emph{IEEE Trans.\ Device
Mater.\ Rel.}, vol.~26, no.~1, pp.~157--166, March 2026.

\end{thebibliography}

\end{document}
